\documentclass[11pt]{article}
\input{../_estilo.tex}
\renewcommand{\tituloejercicio}{Tema 1 --- Ejercicio 7}

\begin{document}

{\LARGE\bfseries\color{azultitulo} Ejercicio 7 --- Mediana de dos arreglos ordenados}

\vspace{2pt}
\textbf{Dificultad:} $\bigstar\bigstar\bigstar$ \quad \textbf{Conferencias:} C2 (El oficio algorítmico),
C3 (Correctitud), C5 (Cotas mínimas)
\vspace{8pt}
\hrule
\vspace{12pt}

\section*{Enunciado}
\begin{enunciado}
Mediana de dos arreglos ordenados de largo $n$: el $n$-ésimo menor de los $2n$ elementos. Diseña un
algoritmo $O(\log n)$, demuestra su correctitud y da la cota de información.
\end{enunciado}

\section{El algoritmo: búsqueda binaria sobre la partición}

Mezclar $A$ y $B$ enteros y mirar la posición $n$ resuelve el problema, pero cuesta $\Theta(n)$; para bajar a
$O(\log n)$ hace falta evitar mirar la mayoría de los $2n$ elementos.

Pensemos en dónde va a quedar el $n$-ésimo menor: a su izquierda (incluyéndolo) tienen que quedar
exactamente $n$ elementos --- la ``mitad baja'' --- repartidos entre $A$ y $B$. Esos dos repartos no son
independientes: si $i$ de esos $n$ elementos vienen de $A$, los otros $j=n-i$ tienen que venir de $B$,
porque juntos suman $n$. Esa es la observación que reduce el problema: en vez de llevar la cuenta por
separado de cuántos elementos de $A$ y cuántos de $B$ quedan a la izquierda, alcanza con un solo número,
$i$ --- el otro, $j$, sale gratis de la relación $i+j=n$. Y una vez fijado $i$ (y por lo tanto $j$), la
respuesta no depende de cómo estén entrelazados los demás elementos: sale de mirar únicamente los cuatro
elementos que quedan justo en la frontera de cada arreglo ($A_{\mathrm{izq}}, A_{\mathrm{der}},
B_{\mathrm{izq}}, B_{\mathrm{der}}$).

Esto convierte el problema de ``encontrar un valor entre $2n$ posibles'' en ``encontrar un índice
$i\in\{0,\dots,n\}$'' --- $n+1$ posibilidades en vez de $2n$. Y ese índice tiene además un comportamiento
particular: a medida que $i$ crece, $A_{\mathrm{izq}}$ (el mayor elemento que $A$ aporta a la mitad baja)
nunca puede bajar, porque $A$ está ordenado; y $B_{\mathrm{der}}$ (el menor elemento que $B$ deja en la
mitad alta) nunca puede subir, porque $j=n-i$ baja y eso corre la frontera de $B$ hacia atrás. Dicho de otra
forma, la condición que separa ``$i$ demasiado chico'' de ``$i$ demasiado grande'' es monótona: cierta de
un lado, falsa del otro. Es una propiedad útil sobre la que se vuelve, con más cuidado, en la sección 2.

\begin{center}
\begin{tikzpicture}[scale=0.62, every node/.style={font=\small}]
  \foreach \c in {0,...,3} { \draw[fill=azulsuave, draw=grisborde] (\c,1) rectangle (\c+1,2); }
  \foreach \c in {4,5} { \draw[fill=white, draw=grisborde] (\c,1) rectangle (\c+1,2); }
  \node[left] at (0,1.5) {$A$:};
  \foreach \c in {0,1} { \draw[fill=azulsuave, draw=grisborde] (\c,-0.5) rectangle (\c+1,0.5); }
  \foreach \c in {2,...,5} { \draw[fill=white, draw=grisborde] (\c,-0.5) rectangle (\c+1,0.5); }
  \node[left] at (0,0) {$B$:};

  \draw[draw=naranjastr, very thick] (3,1) rectangle (4,2);
  \draw[draw=azultitulo, very thick] (4,1) rectangle (5,2);
  \draw[draw=naranjastr, very thick] (1,-0.5) rectangle (2,0.5);
  \draw[draw=azultitulo, very thick] (2,-0.5) rectangle (3,0.5);

  \node[above, naranjastr, font=\scriptsize] at (3.15,2.05) {$A_{\mathrm{izq}}$};
  \node[above, azultitulo, font=\scriptsize] at (4.85,2.05) {$A_{\mathrm{der}}$};
  \node[below, naranjastr, font=\scriptsize] at (1.15,-0.55) {$B_{\mathrm{izq}}$};
  \node[below, azultitulo, font=\scriptsize] at (2.85,-0.55) {$B_{\mathrm{der}}$};

  \draw[decorate,decoration={brace,amplitude=4pt}] (0,2.55) -- (4,2.55);
  \node[above] at (2,2.95) {\scriptsize $i$ elementos};
  \draw[decorate,decoration={brace,amplitude=4pt}] (4,2.55) -- (6,2.55);
  \node[above] at (5,2.95) {\scriptsize $n-i$};

  \draw[decorate,decoration={brace,mirror,amplitude=4pt}] (0,-1.05) -- (2,-1.05);
  \node[below] at (1,-1.45) {\scriptsize $j{=}n{-}i$ elementos};
  \draw[decorate,decoration={brace,mirror,amplitude=4pt}] (2,-1.05) -- (6,-1.05);
  \node[below] at (4,-1.45) {\scriptsize $n-j=i$};
\end{tikzpicture}
\\[14pt]
{\footnotesize División genérica de $A$ y $B$ en una mitad baja (celeste: $i$ elementos de $A$ más
$j=n-i$ de $B$) y una mitad alta (blanco). Las cuatro celdas frontera ---$A_{\mathrm{izq}}$,
$A_{\mathrm{der}}$, $B_{\mathrm{izq}}$, $B_{\mathrm{der}}$--- son las únicas que el algoritmo compara en
cada paso.}
\end{center}

Para ver por qué alcanza con esas cuatro celdas, ayuda mirar qué pasa si de verdad se mezclaran $A$ y $B$
--- cosa que el algoritmo \textbf{nunca hace}, pero que sirve para entender el lema de la sección 2.
\textbf{Ejemplo concreto:} $A=(1,4,6,9)$, $B=(2,3,7,8)$, $n=4$. La partición válida es $i=j=2$, con
$A_{\mathrm{izq}}=4$, $A_{\mathrm{der}}=6$, $B_{\mathrm{izq}}=3$, $B_{\mathrm{der}}=7$.

\begin{center}
\begin{tikzpicture}[scale=0.62, every node/.style={font=\small}]
  \draw[fill=azulsuave, draw=grisborde] (0,1) rectangle (1,2); \node at (0.5,1.5) {1};
  \draw[fill=azulsuave, draw=grisborde] (1,1) rectangle (2,2); \node at (1.5,1.5) {4};
  \draw[fill=white, draw=grisborde] (2,1) rectangle (3,2); \node at (2.5,1.5) {6};
  \draw[fill=white, draw=grisborde] (3,1) rectangle (4,2); \node at (3.5,1.5) {9};
  \draw[draw=naranjastr, very thick] (1,1) rectangle (2,2);
  \draw[draw=azultitulo, very thick] (2,1) rectangle (3,2);
  \node[left] at (0,1.5) {$A$:};

  \draw[fill=azulsuave, draw=grisborde] (0,-0.5) rectangle (1,0.5); \node at (0.5,0) {2};
  \draw[fill=azulsuave, draw=grisborde] (1,-0.5) rectangle (2,0.5); \node at (1.5,0) {3};
  \draw[fill=white, draw=grisborde] (2,-0.5) rectangle (3,0.5); \node at (2.5,0) {7};
  \draw[fill=white, draw=grisborde] (3,-0.5) rectangle (4,0.5); \node at (3.5,0) {8};
  \draw[draw=naranjastr, very thick] (1,-0.5) rectangle (2,0.5);
  \draw[draw=azultitulo, very thick] (2,-0.5) rectangle (3,0.5);
  \node[left] at (0,0) {$B$:};

  \node[font=\scriptsize] at (-0.9,-1.7) {$A\cup B$:};
  % celdas de la mitad baja (celeste), valores 1,2,3,4 con origen A,B,B,A
  \draw[fill=azulsuave, draw=grisborde] (0,-2.1) rectangle (0.5,-1.3); \node[font=\tiny] at (0.25,-1.7) {1};
  \node[font=\tiny, below] at (0.25,-2.15) {A};
  \draw[fill=azulsuave, draw=grisborde] (0.5,-2.1) rectangle (1,-1.3); \node[font=\tiny] at (0.75,-1.7) {2};
  \node[font=\tiny, below] at (0.75,-2.15) {B};
  \draw[fill=azulsuave, draw=grisborde] (1,-2.1) rectangle (1.5,-1.3); \node[font=\tiny] at (1.25,-1.7) {3};
  \node[font=\tiny, below] at (1.25,-2.15) {B};
  \draw[fill=azulsuave, draw=grisborde] (1.5,-2.1) rectangle (2,-1.3); \node[font=\tiny] at (1.75,-1.7) {4};
  \node[font=\tiny, below] at (1.75,-2.15) {A};
  % celdas de la mitad alta (blanco), valores 6,7,8,9 con origen A,B,B,A
  \draw[fill=white, draw=grisborde] (2,-2.1) rectangle (2.5,-1.3); \node[font=\tiny] at (2.25,-1.7) {6};
  \node[font=\tiny, below] at (2.25,-2.15) {A};
  \draw[fill=white, draw=grisborde] (2.5,-2.1) rectangle (3,-1.3); \node[font=\tiny] at (2.75,-1.7) {7};
  \node[font=\tiny, below] at (2.75,-2.15) {B};
  \draw[fill=white, draw=grisborde] (3,-2.1) rectangle (3.5,-1.3); \node[font=\tiny] at (3.25,-1.7) {8};
  \node[font=\tiny, below] at (3.25,-2.15) {B};
  \draw[fill=white, draw=grisborde] (3.5,-2.1) rectangle (4,-1.3); \node[font=\tiny] at (3.75,-1.7) {9};
  \node[font=\tiny, below] at (3.75,-2.15) {A};
  % resaltar los 4 valores frontera dentro de la union mezclada
  \draw[draw=naranjastr, very thick] (1,-2.1) rectangle (1.5,-1.3);
  \draw[draw=naranjastr, very thick] (1.5,-2.1) rectangle (2,-1.3);
  \draw[draw=azultitulo, very thick] (2,-2.1) rectangle (2.5,-1.3);
  \draw[draw=azultitulo, very thick] (2.5,-2.1) rectangle (3,-1.3);
\end{tikzpicture}
\\[4pt]
{\footnotesize La unión $A\cup B$ mezclada y ordenada de verdad (fila de abajo, con la letra de origen de
cada valor). $B_{\mathrm{izq}}{=}3$ y $A_{\mathrm{izq}}{=}4$ son los dos candidatos a ``último de la mitad
baja''; como $4>3$, es $A_{\mathrm{izq}}$ quien queda justo en el corte --- $B_{\mathrm{izq}}$ se queda un
lugar antes. Simétricamente, $A_{\mathrm{der}}{=}6$ gana a $B_{\mathrm{der}}{=}7$ y queda justo después
del corte. Por eso el lema necesita la cadena completa de desigualdades y no alcanza con mirar solo el
elemento que quedó pegado al corte: el otro candidato, aunque no toque el corte, también tiene que estar del
lado correcto.}
\end{center}

\textbf{Especificación (C3 §1).} Precondición: $A$ y $B$ están ordenados de forma creciente y tienen el
mismo largo $n\ge 1$. Postcondición: devuelve el $n$-ésimo menor (1-indexado) de $A\cup B$.

El código real y ejecutable vive en \texttt{verif.py}, en esta carpeta; aquí va el pseudocódigo:

\begin{lstlisting}[style=pseudo]
funcion n_esimo_menor(A, B):
    n (*$\leftarrow$*) longitud(A)
    lo, hi (*$\leftarrow$*) 0, n
    mientras verdadero:
        i (*$\leftarrow$*) (lo+hi) div 2        // i elementos de A en la mitad baja
        j (*$\leftarrow$*) n - i                 // n-i elementos de B en la mitad baja
        Aizq  (*$\leftarrow$*) A[i-1] si i (*$>$*) 0 sino -(*$\infty$*)
        Ader  (*$\leftarrow$*) A[i]   si i (*$<$*) n sino (*$\infty$*)
        Bizq  (*$\leftarrow$*) B[j-1] si j (*$>$*) 0 sino -(*$\infty$*)
        Bder  (*$\leftarrow$*) B[j]   si j (*$<$*) n sino (*$\infty$*)
        si Aizq (*$\le$*) Bder y Bizq (*$\le$*) Ader:      // particion valida
            devolver max(Aizq, Bizq)
        sino si Aizq (*$>$*) Bder:
            hi (*$\leftarrow$*) i - 1                        // tomamos de mas de A
        sino:
            lo (*$\leftarrow$*) i + 1                        // tomamos de menos de A
\end{lstlisting}

\section{Correctitud}

Como en todo argumento de este tipo (C3 §1), correctitud total = correctitud parcial + terminación, y acá
conviene separarlas desde el arranque porque cada una se apoya en una herramienta distinta. La correctitud
parcial responde ``si el ciclo corta, ¿la respuesta es la correcta?'', y se parte a su vez en dos preguntas:
qué significa que una partición sea ``válida'' (el lema del cruce), y por qué existe una partición así para
que el ciclo tenga algo que encontrar (la existencia, por monotonía). La terminación responde, por separado,
``¿el ciclo corta alguna vez?'', y se resuelve con el mismo potencial de cualquier búsqueda binaria: el
rango $[\mathrm{lo},\mathrm{hi}]$ se parte por la mitad en cada vuelta.

\subsection*{Correctitud parcial (I): el lema del cruce}

El algoritmo termina devolviendo $\max(A_{\mathrm{izq}},B_{\mathrm{izq}})$ apenas encuentra una partición que
cumple $A_{\mathrm{izq}}\le B_{\mathrm{der}}$ y $B_{\mathrm{izq}}\le A_{\mathrm{der}}$ --- pero hace falta
justificar por qué esa condición, puntual sobre solo cuatro elementos, garantiza algo sobre los $2n$
elementos enteros. La idea es que esas cuatro desigualdades son exactamente las que hacen falta para que
\textbf{ningún} elemento de la mitad baja (los $i$ de $A$ y $j$ de $B$ elegidos) sea mayor que ningún
elemento de la mitad alta: si eso vale, la mitad baja es forzosamente el conjunto de los $n$ menores de todo
$A\cup B$, sin importar cómo estén entrelazados por dentro.

\begin{resultado}{Lema del cruce}
Si $A_{\mathrm{izq}}\le B_{\mathrm{der}}$ y $B_{\mathrm{izq}}\le A_{\mathrm{der}}$, entonces todo elemento
de $A[0..i{-}1]\cup B[0..j{-}1]$ (la ``mitad baja'', $n$ elementos) es $\le$ todo elemento de
$A[i..n{-}1]\cup B[j..n{-}1]$ (la ``mitad alta'', también $n$ elementos).
\end{resultado}

\textit{Demostración.} Sea $a\in A[0..i{-}1]$ y $b\in B[j..n{-}1]$: $a\le A_{\mathrm{izq}}$ (por ser $A$
creciente) $\le B_{\mathrm{der}}\le b$ (por ser $B$ creciente). Sea $a'\in A[i..n{-}1]$ y
$b'\in B[0..j{-}1]$: simétricamente, $b'\le B_{\mathrm{izq}}\le A_{\mathrm{der}}\le a'$. Las comparaciones
dentro de un mismo arreglo son triviales (está ordenado). $\square$

Como la mitad baja (exactamente $n$ elementos) es $\le$ toda la mitad alta, la mitad baja es \textbf{el
conjunto de los $n$ menores} de $A\cup B$, y su máximo, $\max(A_{\mathrm{izq}},B_{\mathrm{izq}})$, es el
$n$-ésimo menor --- la respuesta que el algoritmo devuelve.

\textbf{Verificado} sobre 2000 pares de arreglos aleatorios, comprobando que, con la partición que encuentra
el algoritmo, el máximo de la mitad baja nunca supera al mínimo de la mitad alta:
\begin{lstlisting}[style=salida]
fallos del lema de cruce (max(izq) <= min(der), deberia ser 0): 0
\end{lstlisting}

\subsection*{Correctitud parcial (II) y terminación: existe una partición válida, y la búsqueda binaria la
encuentra}

El lema del cruce dice qué pasa \emph{si} se encuentra una partición válida; falta mostrar que existe una, y
que el algoritmo efectivamente llega a ella (y no se queda dando vueltas). Como la búsqueda es binaria sobre
$i$, lo que hace falta es exactamente lo que toda búsqueda binaria necesita: que las condiciones de frontera
sean \textbf{monótonas} en $i$.

\textbf{Existencia (monotonía).} Fijemos $i\in\{0,\dots,n\}$ y $j=n-i$, y llamemos $P(i)$ a la condición
``$A_{\mathrm{izq}}\le B_{\mathrm{der}}$'' y $Q(i)$ a ``$B_{\mathrm{izq}}\le A_{\mathrm{der}}$''; una partición
es válida cuando valen las dos. Al aumentar $i$ en 1 (un elemento más de $A$, uno menos de $B$),
$A_{\mathrm{izq}}(i)$ es \textbf{no decreciente} en $i$ ($A$ es creciente), y $B_{\mathrm{der}}(i)=B[n-i]$ es
\textbf{no creciente} en $i$ (porque $j=n-i$ baja). Así que $P$ es \textbf{monótona}: si vale para $i$, vale
para todo $i'<i$; y vale en $i=0$ (ahí $A_{\mathrm{izq}}=-\infty$). Por el mismo argumento con los papeles
cambiados, $Q$ es monótona en la dirección contraria: si vale para $i$, vale para todo $i'>i$; y vale en
$i=n$ (ahí $B_{\mathrm{izq}}=-\infty$). Sea $i_0$ el menor $i$ para el que vale $Q$. Si $i_0=0$, vale también
$P(0)$ y $i_0$ es válida. Si $i_0>0$, $Q(i_0-1)$ es falsa, es decir $B[n-i_0]>A[i_0-1]$, que es exactamente
$P(i_0)$ (pues $A_{\mathrm{izq}}(i_0)=A[i_0-1]$ y $B_{\mathrm{der}}(i_0)=B[n-i_0]$). En los dos casos $i_0$ es
una partición válida: existe al menos una.

\begin{center}
\begin{tikzpicture}[scale=0.7, every node/.style={font=\small}]
  \draw[->] (0,0) -- (6.5,0) node[right, font=\scriptsize] {$i$};
  \draw[->] (0,-0.3) -- (0,3.5) node[above, font=\scriptsize] {valor};
  \draw[naranjastr, thick] plot[smooth] coordinates {(0,0.3) (1,0.6) (2,1.0) (3,1.8) (4,2.6) (5,3.0) (6,3.2)};
  \draw[azultitulo, thick] plot[smooth] coordinates {(0,3.2) (1,3.0) (2,2.6) (3,1.8) (4,1.0) (5,0.6) (6,0.3)};
  \draw[dashed, grisborde] (3,0) -- (3,3.4);
  \node[below, font=\scriptsize] at (3,0) {$i^\ast$};
\node[below, font=\scriptsize] at (0,0) {$0$};
  \node[below, font=\scriptsize] at (6,0) {$n$};
  \node[naranjastr, font=\scriptsize] at (5.3,3.35) {$A_{\mathrm{izq}}(i)$};
  \node[azultitulo, font=\scriptsize] at (5.3,0.05) {$B_{\mathrm{der}}(i)$};
\end{tikzpicture}
\\[4pt]
{\footnotesize $A_{\mathrm{izq}}(i)$ nunca baja, $B_{\mathrm{der}}(i)$ nunca sube. En el caso dibujado se
cruzan en $i^\ast$: para $i<i^\ast$, $P(i)$ (``$A_{\mathrm{izq}}\le B_{\mathrm{der}}$'') es cierta; para
$i>i^\ast$, falsa --- esa es la monotonía que la búsqueda binaria necesita, dibujada sin números concretos.
Con empates puede haber varios índices válidos consecutivos, pero todos devuelven el mismo $n$-ésimo menor.}
\end{center}

\textbf{Terminación.} La búsqueda binaria mueve $\mathrm{hi}=i{-}1$ cuando $A_{\mathrm{izq}}>B_{\mathrm{der}}$
(cuando $i$ es demasiado grande). Si esa condición no se cumple pero la partición aún no es válida, entonces
$B_{\mathrm{izq}}>A_{\mathrm{der}}$ y $i$ es demasiado pequeño, por lo que mueve $\mathrm{lo}=i{+}1$.
Así, cada actualización descarta una mitad del rango de candidatos. Con el potencial $\Phi=\mathrm{hi}-\mathrm{lo}$,
que disminuye estrictamente en cada actualización y queda finalmente en $-1$ cuando el rango se vacía, el
ciclo termina habiendo encontrado una partición válida. Correctitud
parcial (el lema del cruce más la existencia) más terminación dan la \textbf{correctitud total} (C3 §1).

\begin{center}
\begin{tikzpicture}[scale=0.62, every node/.style={font=\small}]
  \node[left] at (-0.6,1) {antes:};
  \foreach \x in {0,...,6} { \filldraw[fill=azulsuave, draw=grisborde] (\x,1) circle (0.18); \node[font=\tiny] at (\x,1.45) {\x}; }
  \filldraw[fill=naranjastr!60, draw=naranjastr] (3,1) circle (0.2);
  \draw[naranjastr, thick] (3.3,0.55) -- (6.3,0.55);
  \node[naranjastr, font=\scriptsize] at (4.8,0.2) {se descartan (monotonía)};

  \node[left] at (-0.6,-0.9) {después:};
  \foreach \x in {0,1,2} { \filldraw[fill=azulsuave, draw=grisborde] (\x,-0.9) circle (0.18); \node[font=\tiny] at (\x,-0.45) {\x}; }
\end{tikzpicture}
\\[4pt]
{\footnotesize Ejemplo con $n=6$: se prueba el candidato del medio, $i=3$ (naranja). Si da ``demasiado
grande'', la monotonía garantiza que $4,5,6$ también lo son --- se descartan sin probarlos, igual que $3$.
Quedan $\{0,1,2\}$: el rango de candidatos para $i$ se parte por la mitad en cada vuelta, de ahí sale
$\Phi=\mathrm{hi}-\mathrm{lo}$ y la cota $O(\log n)$.}
\end{center}

\section{Costo}

Cada vuelta hace $O(1)$ comparaciones y trabajo, y el rango $[\mathrm{lo},\mathrm{hi}]$ (de tamaño $n{+}1$)
se parte por la mitad en cada vuelta --- el mismo potencial de la búsqueda binaria. Costo: $O(\log n)$.

\textbf{Verificación de correctitud} contra fuerza bruta (ordenar y tomar el elemento en la posición
$n{-}1$), sobre arreglos de tamaños $1$ a $5000$:
\begin{lstlisting}[style=salida]
pruebas: 270, fallos (deberia ser 0): 0
\end{lstlisting}

\textbf{Número de iteraciones}, comparado con $\lceil\log_2(n{+}1)\rceil$:
\begin{lstlisting}[style=salida]
     n   iteraciones (peor de 30 pruebas)   techo(log2(n+1))
     1                                  2                  1
     2                                  2                  2
     3                                  3                  2
     5                                  3                  3
    10                                  4                  4
    50                                  6                  6
   200                                  8                  8
  1000                                 10                 10
  5000                                 13                 13
100000                                 17                 17
1000000                                 20                 20
\end{lstlisting}
Para $n$ chico hay una o dos iteraciones de más (overhead de esta implementación concreta, que no está
optimizada al límite); para $n$ grande el crecimiento es $\Theta(\log n)$ con la constante exacta de la
cota, confirmando $O(\log n)$.

\section{Cota de información: \texorpdfstring{$\Omega(\log n)$}{Omega(log n)}}

Para acotar cualquier algoritmo de comparaciones (no solo el de la sección 1), C5 §3 enseña a mirar cuántas
respuestas distintas puede llegar a dar el algoritmo: un árbol de decisión con $h$ niveles tiene a lo sumo
$2^h$ hojas, así que si el problema tiene $k$ respuestas posibles, hace falta $h\ge\log_2 k$. La pregunta
entonces es: ¿cuántas respuestas distintas puede tener este problema?

\begin{resultado}{Observación clave}
La \textbf{identidad} del $n$-ésimo menor --- es decir, cuántos elementos de $A$ hay entre los $n$ menores
--- solo puede tomar $n{+}1$ valores: $i=0,1,\dots,n$.
\end{resultado}

A diferencia de mezclar dos listas completas (C5 §6, que necesita distinguir $\binom{2n}{n}$
entrelazados posibles y cuesta $2n{-}1$), aquí la salida que se pide es mucho más pobre: un solo número
(el valor del $n$-ésimo menor), determinado por un solo parámetro $i\in\{0,\dots,n\}$. Esto sugiere, y
resulta ser, una cota mucho más baja que lineal.

\begin{resultado}{Teorema}
Todo algoritmo de comparaciones necesita, en el peor caso, $\Omega(\log n)$ comparaciones para encontrar el
$n$-ésimo menor de $A\cup B$.
\end{resultado}

\textit{Demostración.} Trabajamos con preguntas binarias del tipo $u\le v$, como las que usa el algoritmo.
El árbol de decisión no enumera las entradas ---hay infinitas---, sino las secuencias de respuestas: cada
nodo interno es una pregunta y cada hoja es una respuesta final. Para forzar que el árbol no pueda ser bajo,
basta exhibir muchas entradas que el algoritmo esté obligado a distinguir entre sí.

Construimos explícitamente $n+1$ instancias con respuestas distintas. Sean
$L_r=-2n+r$ para $0\le r\le n-2$ y $H_r=n+r$ para $1\le r\le n$; los rangos vacíos se omiten.
Para $i=0$, tomamos $A_0=(H_1,\dots,H_n)$ y $B_0=(L_0,\dots,L_{n-2},0)$. Para $1\le i\le n$,
tomamos
\[
A_i=(L_0,\dots,L_{i-2},\,i,\,H_1,\dots,H_{n-i}),\qquad
B_i=(L_{i-1},\dots,L_{n-2},\,H_1,\dots,H_i).
\]
Cada arreglo está ordenado y, en la unión, los primeros $n-1$ valores son los $L_r$, el siguiente es $i$,
y todos los valores $H_r$ son mayores. Por tanto, el $n$-ésimo menor de $A_i\cup B_i$ es exactamente
$i$, y la partición correcta usa exactamente $i$ elementos de $A$.

Fijado un algoritmo arbitrario, estas $n+1$ instancias recorren su árbol de decisión y cada una termina en
alguna hoja. Si dos instancias $A_i,B_i$ y $A_{i'},B_{i'}$ con $i\ne i'$ llevaran a la \textbf{misma hoja},
el algoritmo daría la misma respuesta a las dos --- pero la respuesta correcta depende de $i$, así que al
menos una estaría mal. Como el algoritmo tiene que ser correcto siempre, esto no puede pasar: las $n+1$
instancias necesitan $n+1$ hojas distintas.

Por el Lema 1 de C5 §3, un árbol binario de altura $h$ tiene a lo sumo $2^h$ hojas. Como el árbol necesita
al menos $n+1$ hojas, $2^h\ge n+1$, y despejando, la altura del árbol --- el costo del algoritmo en el peor
caso --- es al menos $h\ge\lceil\log_2(n{+}1)\rceil$. Nótese que $n+1$ y $h$ son cosas distintas: $n+1$ es
cuántas respuestas posibles tiene el \emph{problema}, $h$ es cuántas comparaciones hace el \emph{algoritmo}
--- la desigualdad $2^h\ge n+1$ es el puente que convierte lo primero en una cota sobre lo segundo. Como el
algoritmo era arbitrario, la cota vale para cualquiera. $\square$

\begin{center}
\begin{tikzpicture}[scale=0.8, every node/.style={font=\small}]
  \filldraw[fill=azulsuave, draw=grisborde] (3,3) circle (0.22);
  \filldraw[fill=azulsuave, draw=grisborde] (1.5,2) circle (0.22);
  \filldraw[fill=azulsuave, draw=grisborde] (4.5,2) circle (0.22);
  \draw[grisborde] (3,3) -- (1.5,2);
  \draw[grisborde] (3,3) -- (4.5,2);
  \draw[grisborde] (1.5,2) -- (0.5,1);
  \draw[grisborde] (1.5,2) -- (2.5,1);
  \draw[grisborde] (4.5,2) -- (3.5,1);
  \draw[grisborde] (4.5,2) -- (5.5,1);
  \draw[fill=naranjastr!35, draw=naranjastr] (0.5,0.75) rectangle (1,1.25);
  \draw[fill=naranjastr!35, draw=naranjastr] (2.25,0.75) rectangle (2.75,1.25);
  \draw[fill=naranjastr!35, draw=naranjastr] (3.25,0.75) rectangle (3.75,1.25);
  \draw[fill=naranjastr!35, draw=naranjastr] (5.25,0.75) rectangle (5.75,1.25);
  \node[font=\scriptsize] at (0.75,1) {$i{=}0$};
  \node[font=\scriptsize] at (2.5,1) {$i{=}1$};
  \node[font=\scriptsize] at (3.5,1) {$i{=}2$};
  \node[font=\scriptsize] at (5.5,1) {$i{=}3$};
  \draw[decorate,decoration={brace,amplitude=4pt}] (-0.3,3.15) -- (-0.3,0.6);
  \node[font=\scriptsize, left] at (-0.5,1.9) {$h$ niveles};
  \draw[decorate,decoration={brace,mirror,amplitude=4pt}] (0.3,0.5) -- (5.9,0.5);
  \node[font=\scriptsize, below] at (3.1,0.15) {a lo sumo $2^h$ hojas};
\end{tikzpicture}
\\[10pt]
{\footnotesize Ejemplo con $n=3$ ($n{+}1=4$ instancias): cada instancia de referencia $i=0,1,2,3$ tiene que
terminar en una hoja distinta del árbol de decisión --- si dos cayeran en la misma hoja, el algoritmo daría
la misma respuesta a dos instancias con respuesta correcta distinta. Con $h=2$ niveles alcanzan exactamente
$2^2=4$ hojas, una por instancia; en general, $2^h\ge n+1 \Rightarrow h\ge\log_2(n+1)$.}
\end{center}

El algoritmo de la sección 1 alcanza esta cota hasta una constante, así que es óptimo en el modelo de
comparaciones.

\begin{resultado}{Respuesta a la pregunta del enunciado}
\begin{itemize}[leftmargin=14pt]
\item \textbf{Algoritmo:} búsqueda binaria sobre $i$ (cuántos elementos de $A$ van en la mitad baja), no
sobre los valores; cada paso mira 4 elementos frontera. $O(\log n)$ (secciones 1 y 3).
\item \textbf{Correctitud:} el lema del cruce da la correctitud parcial (partición válida $\Rightarrow$
respuesta correcta); la monotonía de las condiciones de frontera da la existencia de esa partición, y el
potencial $\Phi=\mathrm{hi}-\mathrm{lo}$ da la terminación (sección 2).
\item \textbf{Cota de información:} $\Omega(\log n)$ porque la respuesta solo puede tomar $n+1$ valores
posibles ($i=0,\dots,n$), vía un adversario sobre $n+1$ instancias de referencia y el Lema 1 de C5 §3. El
algoritmo de la sección 1 es óptimo hasta constantes (sección 4).
\end{itemize}
\end{resultado}

\section*{Errores que hay que vigilar}

\begin{alerta}{}
\begin{itemize}[leftmargin=14pt]
\item \textbf{Buscar en el arreglo más largo primero sin ajustar los índices.} Cuando $A$ y $B$ tienen
tamaños distintos (variante más general que la pedida aquí, con $|A|=m\ne n=|B|$), hay que buscar binariamente
sobre el arreglo \textbf{más chico} y calcular $j$ a partir de $i$ y $k$ (la posición buscada) con cuidado
de los límites ($i$ no puede exceder $|A|$ ni hacer que $j$ sea negativo o mayor que $|B|$); mezclar los
roles de $A$ y $B$ a mitad del algoritmo es la fuente más común de errores de implementación de este
problema.
\item \textbf{Pensar que es necesario fusionar (mezclar) las dos listas para encontrar el $n$-ésimo
menor.} Mezclar cuesta $\Theta(n)$ y de hecho \textbf{sí} resuelve el problema, pero es exponencialmente
más lento que necesario: la sección 4 demuestra que $O(\log n)$ basta, porque la pregunta tiene muchas menos
respuestas posibles ($n{+}1$) que la mezcla completa ($\binom{2n}{n}$ entrelazados). Confundir ``resuelve el
problema'' con ``resuelve el problema de forma óptima'' es exactamente la distinción que abre la
conferencia 2 (la quinta pregunta, ¿se puede hacer mejor?, no es retórica).
\item \textbf{Usar $-\infty$/$+\infty$ reales (como \texttt{float('inf')}) en una implementación que
después hace aritmética con esos valores.} Para comparar sirven perfecto; para calcular (por ejemplo, la
mediana promediando dos valores "izquierda" cuando $2n$ es par) hay que tener cuidado de que nunca se
promedie un $\pm\infty$ con un número finito, lo que requiere que $i,j$ estén siempre estrictamente dentro
de rango en el punto final --- una fuente de errores de borde clásica en este problema.
\end{itemize}
\end{alerta}

\end{document}
